\documentclass[../DoAn.tex]{subfiles}

\begin{document}

\begin{center}
    {\bfseries\fontsize{18pt}{22pt}\selectfont TÓM TẮT NỘI DUNG ĐỒ ÁN}
\end{center}
\vspace{1.5cm}

Công tác quản lý ký túc xá tại nhiều trường đại học Việt Nam hiện vẫn dựa chủ
yếu vào quy trình thủ công, tốn nhiều công sức, dễ sai sót và thiếu minh bạch
trong việc phân bổ phòng cũng như giải thích kết quả cho sinh viên. Xuất phát từ
thực trạng đó, đồ án xây dựng \textbf{eDorm} --- một nền tảng tích hợp web và di
động phục vụ quản lý, đăng ký và phân bổ chỗ ở ký túc xá trực tuyến đầu--cuối,
hướng tới mục tiêu tự động hoá, minh bạch hoá và nâng cao trải nghiệm số cho cả
cán bộ quản lý lẫn sinh viên. Về phương pháp, hệ thống thực hiện phân bổ phòng tự
động theo điểm ưu tiên đa tiêu chí, cung cấp môi trường mô phỏng dạng sandbox cho
phép thử nghiệm và khôi phục chỉ bằng một thao tác, đồng thời dự báo nhu cầu chỗ
ở theo bốn kịch bản; phần giao diện được phát triển trên \texttt{React Native}
kết hợp \texttt{Expo} và một Progressive Web App, trong khi không gian cư trú
được trực quan hoá bằng bản đồ ba chiều \texttt{CesiumJS} và tham quan ảo phòng ở
bằng \texttt{Three.js}. Kết quả thực nghiệm cho thấy hệ thống vận hành ổn định ở
mức khoảng 300 người dùng đồng thời và được kiểm thử trên bộ dữ liệu hơn 1.700
sinh viên thuộc bảy toà ký túc xá. Với những kết quả này, \textbf{eDorm} thể hiện
tính khả thi và có thể được triển khai thực tế tại các trường đại học Việt Nam
nhằm hiện đại hoá toàn diện công tác quản lý ký túc xá.

\vspace{2cm}
\begin{flushright}
    Sinh viên thực hiện\\[4pt]
    {\itshape PHAN ĐỨC QUYỀN}
\end{flushright}

\end{document}
